\begin{center}
    \begin{tikzpicture}[show background grid]
        
        \begin{class}[text width=4cm]{Invoker}{0,0}
            \attributett{- command}
            \operationtt{+ setCommand(command)}
            \operationtt{+ executeCommand()}
        \end{class}
        
        \begin{abstractclass}[text width=3cm]{ABCCommand}{5,-0.1}
            \operationtt{+ execute()}
        \end{abstractclass}
        \begin{class}[text width=5.5cm]{CommandA}{2.3,-3}
            \implement{ABCCommand}
            \attributett{- receiver}
            \attributett{- params}
            \attributett{+ CommandA(receiver, params)}
            \attributett{+ execute()}
        \end{class}
        \begin{class}[text width=3cm]{CommandB}{7.7,-3}
            \implement{ABCCommand}
            \operationtt{+ execute()}
        \end{class}
        
        
        
        \begin{class}[text width=3cm]{Client}{-5,-.5}
                    
        \end{class}
        \begin{class}[text width=4cm]{Receiver}{-5,-4}
            \attributett{...}
            \operationtt{+ operation(a, b, c)}
        \end{class}
        
        \draw [ umlcd style , - >] ( Client)  -- ( Receiver);
        \draw [ umlcd style , - >] ( Invoker)  -- ( ABCCommand);
        \draw [ umlcd style dashed line, - >] ( Client)  -- ( CommandA);
    \end{tikzpicture}
\end{center}
